\documentclass[10pt,a4paper]{amsart}
\usepackage[margin=19mm]{geometry}
\usepackage{amsmath,amssymb,mathtools}
\usepackage[scr=rsfs,cal=euler]{mathalfa}
\usepackage{xfrac}
\usepackage[unicode,hidelinks,bookmarks=false]{hyperref}
\newcommand{\ie}{i.e.\ }
\newcommand{\cf}{cf.\ }
\newcommand{\resp}{respectively\ }
\newcommand{\Subsection}{\subsection}
\providecommand{\nobreakdash}{\nobreak\hbox{-}}
\setlength{\emergencystretch}{2em}
\multlinegap0pt
\usepackage{amssymb}
\usepackage{xcolor}
\newtheorem{lemma}{Lemma}
\newtheorem{proposition}{Proposition}
\title{Radicals in primitive axial algebras}
\author{Andrey Mamontov, Sergey Shpectorov, Victor Zhelyabin}
\date{Source: arXiv:2602.11984v1 (CC BY 4.0)}
\begin{document}
\maketitle

\noindent\emph{Source and licence.} Radicals in primitive axial algebras. Version arXiv:2602.11984v1 (CC BY 4.0), distributed by arXiv under the Creative Commons Attribution 4.0 International licence, http://creativecommons.org/licenses/by/4.0/, as recorded in the arXiv licence metadata for this exact version and shown on the abstract page https://arxiv.org/abs/2602.11984v1. Full licence notice: \texttt{licenses/arxiv-2602.11984-CC-BY-4.0.txt}.\par\smallskip

\noindent\textit{Editorial scope.} Counted unit: the article's proposition proposition (source0.tex lines 719--726). The article's complete original proof of it is delivered with it (source0.tex lines 728--761, one \texttt{proof} environment of 34 lines). No mathematics is written, repaired or re-derived: the statement and the proof are the article's own lines, in the article's own notation, and no retained line is substituted or re-flowed.  Retained uncounted as the source-local context the proof needs: lines 330--332; lines 348--356.  Nothing was omitted from any retained span, so the package carries no omission record.  No retained line contains a citation, so the wrapper carries no bibliography and no bibliography entry.  No retained line contains a figure environment, an image-inclusion command or drawing code, so no image is delivered and the package declares no asset dependency.  Editorial formatting only, in addition: an AMS \texttt{amsart} preamble that restates the article's own declarations and macros used by the retained lines, namely the environments \texttt{lemma}, \texttt{proposition} on separate counters, together with the standard packages those lines need.  The retained environments print in this document as Lemma 1, Proposition 1 in order of appearance, whereas the article prints the same items with the numbering of its own full document.  The complete native source of the article as distributed in the arXiv e-print for this exact version is bundled unchanged as \texttt{sources/194511/source0.tex}.
\begin{lemma} \label{block decomposition}
We have that $A=\sum_{a\in X}I_a$.
\end{lemma}

\begin{lemma} \label{sum}
Suppose that an algebra $A$ is decomposed as a sum of its ideals:
%
$$
A=\sum_{j\in J}I_j.
$$
%
If $a\in A$ is a primitive axis then $a\in I_j$ for some $j\in J$.
\end{lemma}

\begin{proposition}
Let $A$ be a primitive axial algebra with $J(A)=0$. Then
%
\begin{enumerate}
\item[(a)] every block $I_c$, $c\in C$, is a simple primitive axial algebra; and
\item[(b)] $A=\oplus_{c\in C}I_c$.
\end{enumerate}
\end{proposition}

\begin{proof}
We first establish (a). Let $I=I_c$ be an arbitrary block. Since $J(A)=0$ and $I\neq 
0$, there exists a maximal ideal $J$ not containing $I$. Then $A=I+J$, since $J$ is 
maximal. Furthermore, $I/I\cap J\cong I+J/J=A/J$ is a simple algebra, which means 
that $I\cap J$ is a maximal ideal of $I$. If $J'$ is any other maximal ideal of $A$ 
then either $I\subseteq J'$ or $I\not\subseteq J'$ and then, similarly, $I\cap J'$ is 
a maximal ideal of $I$. 

We claim that necessarily $I\cap J=I\cap J'$. Suppose that this is not the case, that 
is, $I\cap J\neq I\cap J'$. Then $I=I\cap J+I\cap J'$. By Lemma \ref{sum}, $a\in 
I\cap J$ or $a\in I\cap J'$. In either case, we get a contradiction, because $I$ is 
not the smallest ideal of $A$ containing $a$. Thus, $I\cap J=I\cap J'$.

We have shown that $I\cap J$ is contained in all other maximal ideals $J'$ of $A$. 
Hence, $I\cap J\subseteq J(A)=0$. This gives us that $A=I\oplus J$, since 
$IJ\subseteq I\cap J=0$. However, this means that every ideal of $I$ is an ideal of 
$A$. If $K\neq 0$ is a proper ideal of $I$ then $K+J>J$ is a proper ideal of $A$; a 
contradiction, since $J$ is a maximal ideal. This shows that $I$ has no proper non-
zero ideals, that is, $I$ is simple, proving (a).

For (b), we first note that, for distinct blocks $I_c$ and $I_{c'}$, we have that 
$I_c\cap I_c'=0$, as it is an ideal in both $I_c$ and $I_{c'}$. Hence also 
$I_cI_{c'}\subseteq I_c\cap I_{c'}=0$.

We already know by Lemma \ref{block decomposition} that $A=\sum_{c\in C}I_c$. To see 
that this is a direct sum, suppose we have a non-trivial linear relation 
$\sum_{i=1}^ku_i=0$, where $u_1,u_2,\ldots,u_k$ are elements of pairwise distinct 
blocks $I_{c_1},I_{c_2},\ldots,I_{c_k}$. Assuming that $k$ is smallest possible, we 
must have that each $u_i$ is non-zero. Furthermore, $u_k=-u_1-u_2-\ldots-u_{k-1}\in 
J=I_{c_1}+I_{c_2}+\ldots+I_{c_{k-1}}$. On the one hand, this means that $I_{c_k}\cap 
J\neq 0$ and, since $I_{c_k}$ is simple, this yields $I_{c_k}\subseteq J$. On the 
other hand, $I_{c_k}J=0$ by the previous paragraph. Taking an axis $a\in c_k$, we now 
see that $a=a^2\in I_{c_k}J=0$; this contradiction proves (b).
\end{proof}

\end{document}
